% segment-id: o014.aljabr2.chapter9.duality-rigidity-b.p001
Berdasarkan hubungan timbal balik antara dual kiri dan dual kanan serta Definisi--Proposisi \ref{def:dualizable-obj}, dapat diperiksa bahwa ${}^* (f^*) = f = ({}^* f)^*$; hal ini tampak seketika dengan menggambar ``diagram di dalam diagram''. Selain itu, ${}^* (\identity_X) = \identity_{{}^* X}$ dan $(\identity_X)^* = \identity_{X^*}$ langsung mengikuti definisi.

% segment-id: o014.aljabr2.chapter9.duality-rigidity-b.p002
Berikut dua kelompok identitas berguna mengenai dual kiri/kanan dari $f: X \to Y$; sekali lagi diagram menjadi buktinya.
% segment-id: o014.aljabr2.chapter9.duality-rigidity-b.eq001
\begin{gather}
	\label{eqn:f-dual-coev}
	\begin{tikzpicture}[baseline=(M)]
		\coordinate (A) at (0, 0);
		\coordinate (B) at (1, 0);
		\coordinate (C) at (0, -1);
		\coordinate (D) at (1, -1);
		\draw (C) -- (A) arc[start angle=180, end angle=0, radius=0.5] -- (D);
		\node[draw, fill=white] (M) at ($(A)!.5!(C)$) {$f$};
		\node [below=0.05cm of C] {$Y$};
		\node [below=0.05cm of D] {${}^* X$};
	\end{tikzpicture} = \begin{tikzpicture}[baseline=(M)]
		\coordinate (A) at (0, 0);
		\coordinate (B) at (1, 0);
		\coordinate (C) at (0, -1);
		\coordinate (D) at (1, -1);
		\draw (C) -- (A) arc[start angle=180, end angle=0, radius=0.5] -- (D);
		\node[draw, fill=white] (M) at ($(B)!.5!(D)$) {${}^* f$};
		\node [below=0.05cm of C] {$Y$};
		\node [below=0.05cm of D] {${}^* X$};
	\end{tikzpicture} \qquad \begin{tikzpicture}[baseline=(M)]
		\coordinate (A) at (0, 0);
		\coordinate (B) at (1, 0);
		\coordinate (C) at (0, -1);
		\coordinate (D) at (1, -1);
		\draw (C) -- (A) arc[start angle=180, end angle=0, radius=0.5] -- (D);
		\node[draw, fill=white] (M) at ($(A)!.5!(C)$) {$f^*$};
		\node [below=0.05cm of C] {$X^*$};
		\node [below=0.05cm of D] {$Y$};
	\end{tikzpicture} = \begin{tikzpicture}[baseline=(M)]
		\coordinate (A) at (0, 0);
		\coordinate (B) at (1, 0);
		\coordinate (C) at (0, -1);
		\coordinate (D) at (1, -1);
		\draw (C) -- (A) arc[start angle=180, end angle=0, radius=0.5] -- (D);
		\node[draw, fill=white] (M) at ($(B)!.5!(D)$) {$f$};
		\node [below=0.05cm of C] {$X^*$};
		\node [below=0.05cm of D] {$Y$};
	\end{tikzpicture} \\
% segment-id: o014.aljabr2.chapter9.duality-rigidity-b.eq002
	\label{eqn:f-dual-ev}
	\begin{tikzpicture}[baseline=(M)]
		\coordinate (A) at (0, 0);
		\coordinate (B) at (1, 0);
		\coordinate (C) at (0, -1);
		\coordinate (D) at (1, -1);
		\draw (B) -- (D) arc[start angle=360, end angle=180, radius=0.5] -- (A);
		\node[draw, fill=white] (M) at ($(A)!.5!(C)$) {$f$};
		\node [above=0.05cm of A] {$X$};
		\node [above=0.05cm of B] {$Y^*$};
	\end{tikzpicture} \; = \; \begin{tikzpicture}[baseline=(M)]
		\coordinate (A) at (0, 0);
		\coordinate (B) at (1, 0);
		\coordinate (C) at (0, -1);
		\coordinate (D) at (1, -1);
		\draw (B) -- (D) arc[start angle=360, end angle=180, radius=0.5] -- (A);
		\node[draw, fill=white] (M) at ($(B)!.5!(D)$) {$f^*$};
		\node [above=0.05cm of A] {$X$};
		\node [above=0.05cm of B] {$Y^*$};
	\end{tikzpicture} \qquad \begin{tikzpicture}[baseline=(M)]
		\coordinate (A) at (0, 0);
		\coordinate (B) at (1, 0);
		\coordinate (C) at (0, -1);
		\coordinate (D) at (1, -1);
		\draw (B) -- (D) arc[start angle=360, end angle=180, radius=0.5] -- (A);
		\node[draw, fill=white] (M) at ($(A)!.5!(C)$) {${}^* f$};
		\node [above=0.05cm of A] {${}^* Y$};
		\node [above=0.05cm of B] {$X$};
	\end{tikzpicture} \; = \; \begin{tikzpicture}[baseline=(M)]
		\coordinate (A) at (0, 0);
		\coordinate (B) at (1, 0);
		\coordinate (C) at (0, -1);
		\coordinate (D) at (1, -1);
		\draw (B) -- (D) arc[start angle=360, end angle=180, radius=0.5] -- (A);
		\node[draw, fill=white] (M) at ($(B)!.5!(D)$) {$f$};
		\node [above=0.05cm of A] {${}^* Y$};
		\node [above=0.05cm of B] {$X$};
	\end{tikzpicture}
\end{gather}

% segment-id: o014.aljabr2.chapter9.duality-rigidity-b.prop001
\begin{proposition}
	Misalkan $X \xrightarrow{f} Y \xrightarrow{g} Z$ adalah morfisme dalam kategori monoidal $\mathcal{C}$. Jika semua objek tersebut memiliki dual kiri (atau dual kanan), maka ${}^* (gf) = {}^* f \; {}^* g$ (atau $(gf)^* = f^* g^*$).
\end{proposition}
% segment-id: o014.aljabr2.chapter9.duality-rigidity-b.proof001
\begin{proof}
	Terapkan identitas Zorro \eqref{eqn:Zorro}; satu gambar berbicara lebih banyak daripada seribu kata.
\end{proof}

% segment-id: o014.aljabr2.chapter9.duality-rigidity-b.cor001
\begin{corollary}
	Untuk kategori kaku kiri (atau kaku kanan) $\mathcal{C}$, lengkapi $\mathcal{C}^{\opp}$ dengan struktur monoidal yang membalik arah panah sekaligus urutan variabel $\otimes$. Maka terdapat funktor monoidal $\mathcal{C} \to \mathcal{C}^{\opp}$ yang memetakan objek $X$ ke ${}^* X$ (atau $X^*$) dan morfisme $f$ ke ${}^* f$ (atau $f^*$).
\end{corollary}

% segment-id: o014.aljabr2.chapter9.duality-rigidity-b.p003
Yang mungkin mengejutkan, setiap morfisme antarfunktor monoidal yang berangkat dari kategori kaku kiri atau kanan harus merupakan isomorfisme. Hal ini mencerminkan makna istilah ``kaku''.

% segment-id: o014.aljabr2.chapter9.duality-rigidity-b.prop002
\begin{proposition}\label{prop:monoidal-functor-rigidity}
	Misalkan $F, G: \mathcal{C} \rightrightarrows \mathcal{D}$ funktor monoidal antarkategori monoidal. Jika $\mathcal{C}$ kaku kiri (atau kaku kanan), maka setiap morfisme $\varphi: F \to G$ antarfunktor monoidal merupakan isomorfisme. Bahkan, $(\varphi_X)^{-1} = \left(\varphi_{{}^* X}\right)^*$ (atau ${}^* \left( \varphi_{X^*} \right)$).
\end{proposition}
% segment-id: o014.aljabr2.chapter9.duality-rigidity-b.proof002
\begin{proof}
	Cukup tinjau kasus kaku kiri. Pilih $X \in \Obj(\mathcal{C})$, beserta dual kirinya ${}^* X$ dan data $\mathrm{ev}$, $\mathrm{coev}$. Kita memiliki diagram komutatif
% segment-id: o014.aljabr2.chapter9.duality-rigidity-b.eq003
	\begin{equation*}\begin{gathered}
		\begin{tikzcd}
			\munit \arrow[r, "\sim"] \arrow[rd, "\sim"' sloped] & F(\munit) \arrow[d, "{\varphi_{\munit}}"] \arrow[r, "{F\mathrm{coev}}"] & F(X \otimes {}^* X) \arrow[r, "\sim"] \arrow[d, "{\varphi_{X \otimes {}^* X}}"] & F(X) \otimes F({}^* X) \arrow[d, "{\varphi_X \otimes \varphi_{{}^* X}}"] \\
			& G(\munit) \arrow[r, "{G\mathrm{coev}}"'] & G(X \otimes {}^* X) \arrow[r, "\sim"'] & G(X) \otimes G({}^* X)
		\end{tikzcd} \\
		\begin{tikzcd}
			F({}^* X) \otimes F(X) \arrow[r, "\sim"] \arrow[d, "{\varphi_{{}^* X} \otimes \varphi_X}"'] & F({}^* X \otimes X) \arrow[r, "{F\mathrm{ev}}"] \arrow[d, "{\varphi_{{}^* X \otimes X}}"] & F(\munit) \arrow[r, "\sim"] \arrow[d, "{\varphi_{\munit}}"] & \munit \\
			G({}^* X) \otimes G(X) \arrow[r, "\sim"'] & G({}^* X \otimes X) \arrow[r, "{G\mathrm{ev}}"'] & G(\munit) \arrow[ru, "\sim"' sloped] &
		\end{tikzcd}
	\end{gathered}\end{equation*}
	Selain itu, Proposisi \ref{prop:monoidal-functor-dual} menunjukkan bahwa komposisi sepanjang masing-masing baris diagram membuat $F({}^* X)$ dan $G({}^* X)$ menjadi dual kiri dari $F(X)$ dan $G(X)$. Dengan demikian, $(\varphi_{{}^* X})^*: G(X) \to F(X)$ terdefinisi. Dari sini diperoleh identitas pada diagram

% segment-id: o014.aljabr2.chapter9.duality-rigidity-b.eq004
	\begin{equation*}\begin{aligned}
		\begin{tikzpicture}[baseline=(T), bend angle=70, auto]
			\node (A) {$G(X)$};
			\node (B) [below=of A] {$G(X)$};
			\node (C) [left=of B] {$G({}^* X)$};
			\node (S) [left=of A] {$F({}^* X)$};
			\node[draw, fill=white] (T) at ($(S)!.5!(C)$) {$\varphi_{{}^* X}$};
			\node (D) [left=of S] {$F(X)$}; 
			\node (E) [below=of D] {$G(X)$};
				
			\path (A) edge[->] (B);
			\path (C) edge[bend right] node {$G\mathrm{ev}$} (B);
			\path (C) edge (T);
			\path (T) edge (S);
			\path (D) edge[bend left] node {$F\mathrm{coev}$} (S);
			\path (D) edge[->] (E);
			\node[draw, fill=white] at ($(D)!.5!(E)$) {$\varphi_X$};
		\end{tikzpicture} & = \begin{tikzpicture}[baseline=(A)]
			\node (A) {$G(X)$};
			\node (B) [right=of A] {$G({}^* X)$};
			\node (C) [right=of B] {$G(X)$};
			\path (A) edge[bend left=70] node[midway, auto] {$G(\mathrm{coev})$} (B);
			\path (B) edge[bend right=70] node[midway, auto] {$G(\mathrm{ev})$} (C);	
		\end{tikzpicture}, \\
		\begin{tikzpicture}[baseline=(T), bend angle=70, auto]
			\node (A) {$F(X)$};
			\node (B) [below=of A] {$G(X)$};
			\node (C) [left=of B] {$G({}^* X)$};
			\node (S) [left=of A] {$F({}^* X)$};
			\node[draw, fill=white] (T) at ($(S)!.5!(C)$) {$\varphi_{{}^* X}$};
			\node (D) [left=of S] {$F(X)$}; 
			\node (E) [below=of D] {$F(X)$};
				
			\path (A) edge[->] (B);
			\path (C) edge[bend right] node {$G\mathrm{ev}$} (B);
			\path (C) edge (T);
			\path (T) edge (S);
			\path (D) edge[bend left] node {$F\mathrm{coev}$} (S);
			\path (D) edge[->] (E);
			\node[draw, fill=white] at ($(A)!.5!(B)$) {$\varphi_X$};
		\end{tikzpicture} & = \begin{tikzpicture}[baseline=(A)]
			\node (A) {$F(X)$};
			\node (B) [right=of A] {$F({}^* X)$};
			\node (C) [right=of B] {$F(X)$};
			\path (A) edge[bend left=70] node[midway, auto] {$F(\mathrm{coev})$} (B);
			\path (B) edge[bend right=70] node[midway, auto] {$F(\mathrm{ev})$} (C);	
		\end{tikzpicture},
	\end{aligned}\end{equation*}
	Menurut \eqref{eqn:Zorro}, ruas kanan masing-masing adalah $\identity_{GX}$ dan $\identity_{FX}$. Ini membuktikan bahwa $\left(\varphi_{{}^* X}\right)^*$ memang merupakan invers dari $\varphi_X$.
\end{proof}

% segment-id: o014.aljabr2.chapter9.duality-rigidity-b.prop003
\begin{proposition}\label{prop:dual-internal-Hom}
	Misalkan $\mathcal{C}$ adalah kategori monoidal dan $X, Y \in \Obj(\mathcal{C})$.
	\begin{enumerate}[(i)]
% segment-id: o014.aljabr2.chapter9.duality-rigidity-b.item001
		\item Jika $X$ memiliki dual kiri ${}^* X$, terdapat bijeksi yang saling invers $\Hom_{\mathcal{C}}(X, Y) \xleftrightarrow{1:1} \Hom_{\mathcal{C}}(\munit, Y \otimes {}^* X)$ sebagai berikut.
% segment-id: o014.aljabr2.chapter9.duality-rigidity-b.eq005
		\begin{equation*}\begin{aligned}
			\begin{tikzpicture}[baseline=(M)]
				\node (A) {$X$};
				\node (B) [below=of A]{$Y$};
				\path (A) edge[->] node[midway, draw, fill=white] {$f$} (B);
				\coordinate (M) at ($(A)!.5!(B)$);
			\end{tikzpicture}
			& \; \longmapsto \; \begin{tikzpicture}[baseline=(B), bend angle=70]
				\node (A) {$X$};
				\node (B) [right=of A] {${}^* X$};
				\node (C) [below=of A] {$Y$};
				\node (D) [below=of B] {${}^* X$};
				\path (A) edge[bend left] node[auto] {$\mathrm{coev}_X$} (B);
				\path (A) edge[->] node[midway, draw, fill=white] {$f$} (C);
				\path (B) edge (D);
			\end{tikzpicture}
			\\
			\begin{tikzpicture}[baseline=(D), bend angle=70]
				\node (B) {$X$};
				\node (C) [left=of B] {${}^* X$};
				\node (D) [left=of C] {$Y$};
				\node (A) [below=of D] {$Y$};
				\path (D) edge[->] (A);
				\path (C) edge[bend right] node[auto] {$\mathrm{ev}_X$} (B);
				\path (D) edge[bend left] node[midway, draw, fill=white] {$\varphi$} (C);
			\end{tikzpicture}
			& \; \longmapsfrom \;
			\begin{tikzpicture}[baseline=(A), bend angle=70]
				\node (A) {$Y$};
				\node (B) [right=of A] {${}^* X$};
				\path (A) edge[bend left] node[midway, draw, fill=white] {$\varphi$} (B);
			\end{tikzpicture}
		\end{aligned}\end{equation*}
% segment-id: o014.aljabr2.chapter9.duality-rigidity-b.item002
		\item Jika $X$ memiliki dual kanan $X^*$, terdapat bijeksi yang saling invers $\Hom_{\mathcal{C}}(X, Y) \xleftrightarrow{1:1} \Hom_{\mathcal{C}}(\munit, X^* \otimes Y)$ sebagai berikut.
% segment-id: o014.aljabr2.chapter9.duality-rigidity-b.eq006
		\begin{equation*}\begin{aligned}
			\begin{tikzpicture}[baseline=(M)]
				\node (A) {$X$};
				\node (B) [below=of A]{$Y$};
				\path (A) edge[->] node[midway, draw, fill=white] {$f$} (B);
				\coordinate (M) at ($(A)!.5!(B)$);
			\end{tikzpicture}
			& \; \longmapsto \; \begin{tikzpicture}[baseline=(B), bend angle=70]
				\node (A) {$X$};
				\node (B) [left=of A] {$X^*$};
				\node (C) [below=of A] {$Y$};
				\node (D) [below=of B] {$X^*$};
				\path (B) edge[bend left] node[auto] {$\mathrm{coev}_X$} (A);
				\path (A) edge[->] node[midway, draw, fill=white] {$f$} (C);
				\path (B) edge (D);
			\end{tikzpicture}
			\\
			\begin{tikzpicture}[baseline=(D), bend angle=70]
				\node (B) {$X$};
				\node (C) [right=of B] {$X^*$};
				\node (D) [right=of C] {$Y$};
				\node (A) [below=of D] {$Y$};
				\path (D) edge[->] (A);
				\path (B) edge[bend right] node[auto] {$\mathrm{ev}_X$} (C);
				\path (C) edge[bend left] node[midway, draw, fill=white] {$\varphi$} (D);
			\end{tikzpicture}
			& \; \longmapsfrom \;
			\begin{tikzpicture}[baseline=(A), bend angle=70]
				\node (A) {$X^*$};
				\node (B) [right=of A] {$Y$};
				\path (A) edge[bend left] node[midway, draw, fill=white] {$\varphi$} (B);
			\end{tikzpicture}
		\end{aligned}\end{equation*}
% segment-id: o014.aljabr2.chapter9.duality-rigidity-b.item003
		\item Lebih umum lagi, apabila dual-dual yang dimaksud ada, terdapat bijeksi kanonik
% segment-id: o014.aljabr2.chapter9.duality-rigidity-b.eq007
		\[\begin{tikzcd}[row sep=tiny]
			\Hom_{\mathcal{C}}\left( W \otimes X, Y \right) \arrow[r, "1:1"] & \Hom_{\mathcal{C}}(W, Y \otimes {}^* X) \\
			f \arrow[mapsto, r] & (f \otimes \identity_{{}^* X}) (\identity_W \otimes \mathrm{coev}_X)
		\end{tikzcd}\]
		dan
% segment-id: o014.aljabr2.chapter9.duality-rigidity-b.eq008
		\[\begin{tikzcd}
			\Hom_{\mathcal{C}}\left( X \otimes W, Y \right) \arrow[r, "1:1"] & \Hom_{\mathcal{C}}(W, X^* \otimes Y) \\
			f \arrow[mapsto, r] & (\identity_{X^*} \otimes f) (\mathrm{coev}_X \otimes \identity_W).
		\end{tikzcd}\]
	\end{enumerate}
\end{proposition}
% segment-id: o014.aljabr2.chapter9.duality-rigidity-b.proof003
\begin{proof}
	Untuk (i) dan (ii), pemeriksaan bahwa pemetaan-pemetaan tersebut saling invers hanyalah penerapan \eqref{eqn:Zorro}. Dengan gagasan yang sama, suatu diagram yang sedikit lebih rumit untuk digambar cukup untuk menjelaskan (iii).
\end{proof}

% segment-id: o014.aljabr2.chapter9.duality-rigidity-b.p004
Dalam konteks (i) dan (ii), mudah mendeskripsikan komposisi morfisme pada unsur $\varphi$ di ruas kanan bijeksi (petunjuk: terapkan morfisme $\mathrm{ev}$). Operasi-operasi kanonik ini menghasilkan pernyataan berikut.

% segment-id: o014.aljabr2.chapter9.duality-rigidity-b.cor002
\begin{corollary}\label{prop:rigid-Hom-internal}
	\index[sym1]{Hom-internal}
	Misalkan $\mathcal{C}$ adalah kategori kaku kiri (atau kaku kanan) sebagaimana dalam Definisi \ref{def:rigid-cat}. Dengan mendefinisikan $\Hom$ internal sebagai $\iHom_{\text{kiri}}(X, Y) := Y \otimes {}^* X$ (atau $\iHom_{\text{kanan}}(X, Y) := X^* \otimes Y$), kategori $\mathcal{C}$ menjadi kategori monoidal tertutup kanan (atau tertutup kiri) sebagaimana dalam Definisi \ref{def:closed-monoidal-cat}.
\end{corollary}
% segment-id: o014.aljabr2.chapter9.duality-rigidity-b.proof004
\begin{proof}
	Terapkan Proposisi \ref{prop:dual-internal-Hom} (iii) pada Definisi \ref{def:closed-monoidal-cat}; pemeriksaan selebihnya bersifat rutin.
\end{proof}

% segment-id: o014.aljabr2.chapter9.duality-rigidity-b.cor003
\begin{corollary}\label{prop:rigid-exact}
	Misalkan $X \in \Obj(\mathcal{C})$.
	\begin{itemize}
% segment-id: o014.aljabr2.chapter9.duality-rigidity-b.item004
		\item Jika $X$ memiliki dual kiri ${}^* X$, maka $(\cdot) \otimes X$ mempertahankan semua $\varinjlim$ kecil, sedangkan $X \otimes (\cdot)$ mempertahankan semua $\varprojlim$ kecil.
% segment-id: o014.aljabr2.chapter9.duality-rigidity-b.item005
		\item Jika $X$ memiliki dual kanan $X^*$, maka $(\cdot) \otimes X$ mempertahankan semua $\varprojlim$ kecil, sedangkan $X \otimes (\cdot)$ mempertahankan semua $\varinjlim$ kecil.
	\end{itemize}
\end{corollary}
% segment-id: o014.aljabr2.chapter9.duality-rigidity-b.proof005
\begin{proof}
	Tinjau kasus ketika $X$ memiliki dual kiri ${}^* X$. Proposisi \ref{prop:dual-internal-Hom} (iii) menyiratkan bahwa funktor $(\cdot) \otimes X$ memiliki adjoin kanan $(\cdot) \otimes {}^* X$, sedangkan funktor $X \otimes (\cdot) \simeq ({}^* X)^* \otimes (\cdot)$ memiliki adjoin kiri ${}^* X \otimes (\cdot)$.
	
	Dengan menukar urutan $\otimes$ dalam argumen di atas, kita memperoleh kasus ketika $X$ memiliki dual kanan $X^*$.
\end{proof}
